\appendix
\numberwithin{equation}{section}
\numberwithin{table}{section}

\section{Detailed model and training equations}
\label{app:detailed-equations}

This appendix collects the parameterization and training definitions omitted
from the main narrative.  The redshift variables relative to the fixed pivot
$z_*=0.5$ are
\begin{equation}
  r(z)=\frac{1+z}{1+z_*},
  \qquad
  r_q=\frac{1+z_q}{1+z_*}.
\end{equation}
The broad redshift trend and smooth luminosity-like transition are
\begin{equation}
  \begin{aligned}
    R_z(z) &= r(z)^\eta,\\
    R_L(z) &=
    \left[
      \frac{1+[r(z)/r_q]^s}{1+(1/r_q)^s}
    \right]^{\xi/s}.
  \end{aligned}
\end{equation}

The characteristic scale and shape quantities evolve as
\begin{equation}
  \begin{aligned}
    k_t(z) &= k_{t,*}\,r(z)^{\gamma_t}, &
    n(z) &= n_*\,r(z)^{\gamma_n},\\
    \alpha(z) &= \alpha\,r(z)^{\gamma_\alpha}, &
    m(z) &= m\,r(z)^{\gamma_m}.
  \end{aligned}
\end{equation}
With $x=k/k_t(z)$, the scale response is
\begin{equation}
  S(k,z)=1+q\,
  \frac{x^{n(z)}}{1+x^{n(z)}}
  \left(\frac{1+x^{m(z)}}{2}\right)^{\alpha(z)/m(z)}.
\end{equation}
Together these factors give the positive surface in Equation~(8) of the main
text.

For the autoencoder, the pointwise logarithm and training-only normalization
are
\begin{equation}
  \begin{aligned}
    Y(k,z) &= \log_{10}\mathcal A_\Theta(k,z),\\
    \bm{x} &= \frac{Y-\overline{Y}}{\sigma}.
  \end{aligned}
\end{equation}
The encoder and decoder maps are
\begin{equation}
  \bm{z}=f_{\phi}(\bm{x}),
  \qquad
  \widehat{\bm{x}}=g_{\theta}(\bm{z}),
  \qquad
  \bm{z}\in\mathbb{R}^{2}.
\end{equation}
Finally, the training loss is the mean squared reconstruction error
\begin{equation}
  \mathcal L_{\rm MSE}
  =\frac{1}{N N_z N_k}
  \sum_{i=1}^{N}\sum_{j=1}^{N_z}\sum_{\ell=1}^{N_k}
  \left(\widehat{x}_{ij\ell}-x_{ij\ell}\right)^2,
\end{equation}
where $N$ is the batch size, $N_z=31$, and $N_k=101$.

\subsection{Sampling bounds}

Table~\ref{tab:sampling-bounds} records the data-generation bounds used to
construct the synthetic response family.  They are design choices for this
compression experiment, not observational constraints or posterior
distributions from a galaxy survey.

\begingroup
\small
\renewcommand{\arraystretch}{1.05}
\begin{longtable}{@{}p{0.15\linewidth}p{0.32\linewidth}p{0.19\linewidth}p{0.21\linewidth}@{}}
\caption{Sampling bounds for the 13-parameter synthetic model family.}
\label{tab:sampling-bounds}\\
\toprule
Parameter & Main effect & Range & Sampling \\
\midrule
\endfirsthead
\toprule
Parameter & Main effect & Range & Sampling \\
\midrule
\endhead
$\eta$ & broad redshift tilt & $[-1.5,1.5]$ & uniform \\
$\xi$ & strength of second redshift trend & $[-1.5,1.5]$ & uniform \\
$s$ & sharpness of its transition & $[1,10]$ & uniform \\
$z_q$ & transition redshift & $[0.5,2.5]$ & uniform \\
$q$ & size of scale response & $[0,2]$ & uniform \\
$n_*$ & scale-transition sharpness at the pivot & $[1,5]$ & uniform \\
$k_{t,*}$ & pivot transition wavenumber & $[0.1,1]$ & log-uniform \\
$\alpha$ & high-$k$ slope at the pivot & $[-0.4,0.4]$ & uniform \\
$m$ & smoothness of the high-$k$ tail at the pivot & $[1,4]$ & uniform \\
$\gamma_t$ & redshift evolution of $k_t$ & $[-0.5,0.5]$ & uniform \\
$\gamma_n$ & redshift evolution of $n$ & $[-0.35,0.35]$ & uniform \\
$\gamma_\alpha$ & redshift evolution of $\alpha$ & $[-0.25,0.25]$ & uniform \\
$\gamma_m$ & redshift evolution of $m$ & $[-0.2,0.2]$ & uniform \\
\bottomrule
\end{longtable}
\endgroup

\subsection{Complete training configuration}

Table~\ref{tab:training-config-details} preserves the numerical implementation
details omitted from the shorter main-text table.  The values come from
\path{Config/NLA/Training/Standard.yaml} and
\path{Config/NLA/Surface/Standard.yaml}.

\begingroup
\small
\renewcommand{\arraystretch}{1.08}
\begin{longtable}{@{}p{0.29\linewidth}p{0.62\linewidth}@{}}
\caption{Complete settings for the shared autoencoder training template.}
\label{tab:training-config-details}\\
\toprule
Setting & Value \\
\midrule
\endfirsthead
\toprule
Setting & Value \\
\midrule
\endhead
Training objective & MSE in normalized $\log_{10}\mathcal A_\Theta$ space \\
Optimization & AdamW; learning rate $2\times10^{-4}$; weight decay $10^{-5}$ \\
Batching & training batch 512; evaluation batch 512 \\
Numerical controls & gradient-norm limit 1; mixed precision enabled \\
Learning-rate schedule & reduce on plateau; factor $0.5$; patience 12 epochs;
minimum learning rate $10^{-7}$ \\
Termination & maximum 1,000 epochs; early-stopping patience 50 epochs;
minimum improvement $10^{-7}$ \\
Checkpoint criterion & lowest complete-validation MSE \\
Accuracy target & recovered variance at least $0.999$ \\
Randomness and device & seed 42; deterministic kernels not enforced;
device selected automatically \\
\bottomrule
\end{longtable}
\endgroup

\section{Reproducibility checklist}
\label{app:reproducibility}

\begin{itemize}
  \item Code version: Git commit \texttt{fb975ed} in the IAFlowCloud
        repository.
  \item Authoritative dataset: \texttt{Data/NLA/Samples/NLA.hdf5}, containing
        coordinates, the 13 shape parameters, factorized components, and the
        stored train/validation/test indices.
  \item Parameter order and sampling bounds:
        Appendix Table~\ref{tab:sampling-bounds}.
  \item Coordinate grid and split sizes: 31 redshifts from $z=0$ to $z=3$,
        101 logarithmically spaced wavenumbers from
        $k=0.01$ to $10\,{\rm Mpc}^{-1}$, and an $80{,}000/10{,}000/10{,}000$
        split (Section~\ref{sec:model}).
  \item PCA artifacts: \texttt{Data/NLA/PCA/}, including
        \texttt{PCAValidationMetrics.json} and the rank-25 joblib transform.
        The physical-model and PCA notebooks use the CosmoConda environment
        (Python~3.12, NumPy~2.0, scikit-learn~1.7, PyCCL~3.2).
  \item Reported autoencoder: Conv1D Depth03 with latent dimension 2,
        2,039,969 parameters, 25 training epochs, best checkpoint at epoch 20.
        The later shared training template in
        \texttt{Config/NLA/Training/Standard.yaml} uses seed 42 and a 1,000-epoch
        maximum; that longer schedule is not the source of
        Table~\ref{tab:preliminary-results}.
  \item Random seeds: dataset split seed stored in the HDF5 generation
        metadata; reported autoencoder seed 42 in the shared training
        template.  Only one short autoencoder run is claimed here.
  \item Final test evaluation: not run.  The complete test split is reserved
        until architecture and training choices are frozen.
\end{itemize}

\section{Implementation notes}

The physical model lives in the \texttt{ia\_models.nla} package.  The learning
workflow lives in \texttt{iaflow}.  Shared surface data are cached once as
$\log_{10}\mathcal A_\Theta$ with training-only mean and global RMS
normalization.  Evaluation, diagnostics, and latent export read a run's
resolved configuration rather than an edited YAML file, so a checkpoint cannot
silently change its architecture after training.
